\section{What Daylight Does---and Does Not---Mean Astronomically}

The textual claim needs almost no astronomy. Stars other than the sun are
present above the horizon in daylight; atmospheric scattering normally prevents
the unaided eye from distinguishing them. Bright planets and, under exceptional
conditions, a very small number of stellar objects may be visible, but that
does not turn the host of heaven into a countable daylight display. The
chapter's challenge is not solved by observing that stars continue to exist
when unseen.

Nor should modern optical possibilities be projected into the scene. A deep
well, eclipse, telescope, or extraordinary atmospheric condition can alter
visibility, but Genesis names none of them.

The safer conclusion is deliberately modest. On the continuous reading, the
command's ordinary visual object is unavailable in the ordinary way. Whether
that unavailability is the intended point cannot be established from astronomy.
The narrative may be visionary, may omit an interval, or may presuppose divine
enablement. Physical possibility sets limits on a proposed scene; it does not
select among the exegetical answers.

